\begin{table}[ht]\centering\small
\begin{tabular}{rrrrrrrr}
\toprule
& & \multicolumn{2}{c}{$\eta=0.3$} & \multicolumn{2}{c}{$\eta=0.2$} & \multicolumn{2}{c}{$\eta=0.1$}\\
$L$ & $M$ & 격자 & LLL & 격자 & LLL & 격자 & LLL\\
\midrule
31 & 5 & 1.78 & 1.79 & 2.01 & 2.02 & 2.02 & 2.02\\
43 & 6 & 2.04 & 2.04 & 2.58 & 2.58 & 2.83 & 2.83\\
53 & 7 & 2.56 & 2.56 & 2.56 & 2.56 & 3.52 & 3.52\\
61 & 8 & 2.56 & 2.56 & 2.56 & 2.56 & 4.50 & 4.50\\
71 & 9 & 2.73 & 2.73 & 3.45 & 3.45 & 3.90 & 3.90\\
73 & 10 & 3.11 & 3.11 & 3.58 & 3.58 & 3.90 & 3.90\\
101 & 11 & 2.81 & 2.81 & 3.18 & 3.18 & 5.59 & 5.59\\
103 & 12 & 2.81 & 2.81 & 4.23 & 4.23 & 5.85 & 5.85\\
109 & 13 & 2.81 & 2.81 & 4.61 & 4.61 & 6.64 & 6.94\\
113 & 14 & 3.73 & 3.73 & 4.94 & 4.94 & 7.04 & 7.04\\
139 & 15 & 4.41 & 4.41 & 5.11 & 5.11 & 7.86 & 7.86\\
157 & 16 & 4.24 & 4.24 & 5.38 & 5.38 & 7.32 & 7.32\\
173 & 17 & 3.42 & 5.22 & 6.30 & 6.52 & 8.87 & 9.28\\
181 & 18 & 3.42 & 3.42 & 6.34 & 6.34 & $>9.10$ & 10.77\\
191 & 19 & 4.87 & 4.87 & 6.81 & 7.83 & $>9.08$ & 10.60\\
193 & 20 & 4.78 & 4.78 & 6.81 & 6.81 & 9.14 & 11.90\\
199 & 21 & 4.78 & 4.78 & 6.81 & 8.73 & $>9.08$ & 12.70\\
239 & 22 & 5.79 & 6.16 & 8.41 & 8.41 & $>9.08$ & 11.19\\
241 & 23 & 6.20 & 6.79 & 8.55 & 9.67 & $>9.05$ & 11.77\\
251 & 24 & 6.15 & -- & 8.55 & -- & $>9.15$ & --\\
269 & 25 & 5.75 & -- & 7.35 & -- & $>9.07$ & --\\
271 & 26 & 5.75 & -- & 7.35 & -- & $>9.02$ & --\\
283 & 27 & 5.75 & -- & $>8.99$ & -- & $>8.99$ & --\\
\bottomrule
\end{tabular}
\caption{첫 공명 높이 $\log_{10}\tau^*_{10}(L,\eta)$. ``격자'' 열: 인증된 값 (\texttt{code/resscan.c}; 정밀도 $10^{-3}$ 이하) 또는
$L$ 당 120\,s 안에 공명이 없을 때의 인증된 하한 ``$>$''. ``LLL'' 열: \texttt{code/lll\_res.py} 가 찾은 점 $t$ ($g_L(t)\ge(1-\eta)M$ 을 mpmath 로 확인) 의
$\log_{10}t$ --- 즉 $\tau^*_{10}\le t$ 의 상한. $L$ 은 각 $M$ 을 처음 달성하는 소수. \NUMERIC\
(\texttt{out/resscan\_A.txt}, \texttt{out/resscan\_B.txt}, \texttt{out/lll\_res.txt}, \texttt{data/resonance.csv}, \texttt{code/res\_analyze.py})}\label{NUM:tab-res}
\end{table}

\begin{proposition}[공명 높이 자료]\label{NUM:prop-res}
\NUMERIC\ (\texttt{code/resscan.c}, \texttt{code/lll\_res.py}, \texttt{code/resolve\_amb.py}, \texttt{code/res\_analyze.py} $\to$ \texttt{out/res\_fit.txt}, \texttt{data/resonance.csv})
\begin{enumerate}[label=(\arabic*)]
\item 표~\ref{NUM:tab-res} 의 모든 ``격자'' 항목은 보조정리~\ref{NUM:lem-grid} 의 인증서로 얻었다. 판정 불가 후보는 한 번 ($L=193$, $\eta=0.3$, $t\approx60322.49$) 나왔고,
국소 최대화 (\texttt{out/resolve\_amb.txt}) 로 그 칸에는 공명이 없고 첫 교차가 $t_c=60322.5932$ 임을 확인하였다 (\texttt{resscan} 의 보고값과 일치).
\item LLL 이 찾은 점은 인증된 공명 없는 구간 안에 한 번도 들어가지 않았다 (모순 없음). 두 값이 모두 있는 52 경우 중 42 경우에서
LLL 점은 격자의 첫 교차점과 $3$ 이내 (같은 공명 봉우리) 이다. 즉 이 범위에서 LLL 은 사실상 \emph{첫} 공명을 찾는다.
\item 최소제곱 기울기 (상용로그, 소수 하나당):
\begin{itemize}
\item $\eta=0.3$: 격자 $\tau^*$ (23개, $M\in[5,27]$) $\log_{10}\tau^*\approx0.83+0.199M$; LLL 점 (19개, $M\le23$) $\log_{10}t\approx0.57+0.229M$; 무작위 위상 모형 기울기 0.250
\item $\eta=0.2$: 격자 $\tau^*$ (22개, $M\in[5,26]$) $\log_{10}\tau^*\approx0.45+0.315M$; LLL 점 (19개, $M\le23$) $\log_{10}t\approx-0.46+0.403M$; 무작위 위상 모형 기울기 0.354
\item $\eta=0.1$: 격자 $\tau^*$ (14개, $M\in[5,20]$) $\log_{10}\tau^*\approx-0.13+0.496M$; LLL 점 (19개, $M\le23$) $\log_{10}t\approx-0.99+0.595M$; 무작위 위상 모형 기울기 0.520
\end{itemize}
\item $\eta=0.1$ 에서 격자는 $M\ge18$ 부터 120\,s 안에 공명을 찾지 못한다: 인증된 하한 $\tau^*_{10}>9.81e+08$ (최소값).
\end{enumerate}
\end{proposition}

\begin{remark}[해석]\label{NUM:rem-res-heur}
\HEUR\ (\texttt{code/ratefn.py} $\to$ \texttt{out/ratefn.txt})
(i) 위상 $t\log p$ ($p\in P$) 를 독립 균등으로 보는 무작위 위상 모형에서는 $\Pr(|\frac1M\sum\e^{i\theta_p}|\ge x)=\e^{-M I(x)+O(\log M)}$,
$I(x)=\sup_\lambda(\lambda x-\log I_0(\lambda))$ (Cramér; 2차원 벡터 $(\cos\theta,\sin\theta)$ 의 로그적률함수가 $\log I_0(|\lambda|)$), 상관 길이는 $O(1)$ 이므로
$\log\tau^*\approx M I(1-\eta)$ 가 예측된다: $I(0.7)=0.576$, $I(0.8)=0.815$, $I(0.9)=1.196$ (상용로그 기울기 $0.250,0.354,0.520$).
관측 기울기는 $\eta=0.2,0.1$ 에서 이 예측과 가깝고, $\eta=0.3$ 에서는 작다 (작은 $M$ 에서 공유 공명 --- 예: $\eta=0.3$ 에서 $L=53,61$ 은 $t\approx360$, $L=101,103,109$ 는 $t\approx642$ 에서 공명 --- 과 유한 크기 효과).
(ii) Dirichlet 동시근사는 $\tau^*\lesssim Q^M$, $Q=\lceil2\pi/\arccos(1-\eta)\rceil=8,10,14$ ($\log_{10}Q=0.90,1.00,1.15$) 를 주는데, 관측 기울기는 그 $1/5$--$1/2$ 로
Dirichlet 상한은 지수에서 날카롭지 않다. DESIGN S4 의 ``$T\ge(C/\eta)^M$ 이면 NR 실패'' 는 지수 높이가 자연스러운 한계라는 점에서 자료와 부합한다.
(iii) VKGAP-C 에 대한 함의: 자료는 $\log\tau^*_{10}(L,\eta)\asymp M$ (선형 증가) 과 부합하고, 따라서 $\mathrm{NR}(P,\e^{cM},\eta)$ 는 $c<c(\eta)$ 로
기대할 수 있다 ($c(0.1)\approx0.5\ln10\approx1.1$ 정도; 이 범위 $M\le33$ 에서의 추정일 뿐). 단 명제~\ref{NUM:prop-t1} 에 의해 절단은 $|\tau|\ge1$ 이 아니라
$\eta<0.0194$ 또는 $|\tau|\ge\tau_0(\eta)$ 여야 한다. $\eta=0.01$ 의 높이는 $\tau^*(0.1)$ 보다 크므로 (주의~\ref{NUM:rem-t1}) 위 하한은 $\eta=0.01$ 에도 유효하다.
이 모든 것은 $M\le33$, $\ln L\le6$ 의 자료에 근거한 외삽이며 증명이 아니다.
\end{remark}
